% !PRINTED 173
\weberpage{173}
Bedeutet $A$ eine Substitution der endlichen Gruppe $S$, so
können wir aus $\Theta(x)$ eine neue Function:
\weq{2}{\Theta[A(x)] = c_1' x_1 + c_2' x_2 + \cdots + c_n x_n} % [sic]
ableiten, deren Co\"efficienten $c_i'$ nach \S.\,37, 9. mit den ursprünglichen
Co\"efficienten $c_i$ durch die zu $A$ \gesperrt{transponirte Substitution}
\[ (c') = A_1(c) \]
zusammenhängen.

Wenn nun $A, B, C, \dots$ die sämmtlichen Substitutionen der
Gruppe $S$ sind, so kann man in gleicher Weise die Functionen
\weq{3}{\Theta[A(x)],\ \Theta[B(x)],\ \Theta[C(x)],\ \dots}
bilden, die wir auch kürzer durch
\weq{4}{\Theta,\ \Theta_1,\ \Theta_2,\ \dots,\ \Theta_{\mu-1}}
bezeichnen, wenn $\mu$ der Grad der Gruppe $S$ ist.

Nun kann man über die Co\"efficienten $c_i$, die bis jetzt noch
ganz willkürlich sind, so verfügen, dass keine zwei der Functionen
(4) mit einander identisch werden (Bd.\ I, \S.\,143, 1.).
Aus (3) aber ergiebt sich:

\gesperrt{1. Wenn man in den $\mu$ Functionen (4) gleichzeitig
irgend eine Substitution aus $S$ anwendet, so
ändert sich die Gesammtheit dieser Functionen
nicht, sondern sie erleiden nur eine Permutation.}

Daraus ergiebt sich ferner:

\gesperrt{2. Jede symmetrische Function der Grössen (4) ist
eine absolute Invariante der Gruppe $S$.}

Wir bemerken noch, dass man an Stelle der Function $\Theta$
irgend eine andere auch nicht lineare, selbst eine gebrochene
oder inhomogene Function setzen könnte, wenn nur die $\mu$ Functionen
(4) von einander verschieden sind.

Bilden wir nun das Product:
\weq{5}{\begin{aligned}
\Phi(t) &= (t-\Theta)(t-\Theta_1)\dots(t-\Theta_{\mu-1}) \\
&= t^\mu + A_1 t^{\mu-1} + A_2 t^{\mu-2} + \cdots + A_\mu,
\end{aligned}}
welches eine ganze rationale Function $\mu^{\text{ten}}$ Grades von $t$ ist, so
sind die Co\"efficienten $A_1, A_2, \dots, A_\mu$ nach 2. Invarianten der
Gruppe $S$, und die Grössen $\Theta, \Theta_1, \dots, \Theta_{\mu-1}$ sind die Wurzeln
der Gleichung
\weq{6}{\Phi(t) = 0.}
